\chapter{INTRODUCTION}

\section{Overview}

Food logging became the starting point for this project because the manual process is easy to abandon. A user has to remember what was eaten, estimate the quantity, search for a matching food item, and repeat the same steps for every side item. This is especially awkward for Indian meals. A single plate may contain rice, dal, chutney, curry, fried bread, pickle, and a sweet, sometimes in small portions that are still nutritionally relevant. The meal is obvious to the person eating it, but a logging application usually treats it as several separate search tasks.

Computer vision can reduce the entry work, but it cannot be treated as the whole solution. Food changes shape after cooking, lighting changes colour, and bowls or plates can hide parts of the item. Indian dishes also share ingredients and visual cues. In some samples, \code{aloo_gobi}, \code{dum_aloo}, and \code{aloo_masala} were visually close enough that the prediction had to be checked against the confusion matrix and sample images. Chutneys caused another kind of problem because they often appeared as small side portions rather than main objects. These observations made the project more than a detector-training exercise.

The implemented prototype therefore connects recognition with nutrition lookup. It uses a YOLO object detector, a curated Indian food dataset, a SQLite nutrition database derived mainly from the Indian Nutrient Database, a FastAPI backend, and a Next.js frontend. After an image is uploaded, the backend detects visible food items, maps each detected class to a nutrition entry, calculates macronutrients, and prepares suggestions using the user's goal and selected health context.

\begin{figure}[H]
\centering
\resizebox{0.96\textwidth}{!}{%
\begin{tikzpicture}[
    box/.style={draw, rounded corners, align=center, minimum width=3.3cm, minimum height=1.05cm, fill=blue!6},
    data/.style={draw, rounded corners, align=center, minimum width=3.3cm, minimum height=1.05cm, fill=green!8},
    outputbox/.style={draw, rounded corners, align=center, minimum width=3.3cm, minimum height=1.05cm, fill=orange!10},
    arrow/.style={-Latex, thick},
    node distance=0.8cm and 0.9cm
]
\node[box] (upload) {Meal image\\uploaded by user};
\node[box, right=of upload] (detect) {YOLO detects\\visible food items};
\node[data, right=of detect] (map) {Class labels map\\to nutrition rows};
\node[outputbox, right=of map] (result) {Macros and\\diet suggestions};
\node[data, below=of detect] (dataset) {72-class curated\\Indian food dataset};
\node[data, below=of map] (db) {SQLite nutrition\\database and mappings};
\draw[arrow] (upload) -- (detect);
\draw[arrow] (detect) -- (map);
\draw[arrow] (map) -- (result);
\draw[arrow] (dataset) -- (detect);
\draw[arrow] (db) -- (map);
\end{tikzpicture}
}
\caption{End-to-end workflow of the implemented food recognition and nutrition system}
\label{fig:project-workflow}
\end{figure}

\section{Motivation}

The motivation came from two issues seen during development. The first was speed of use. If the user has to spend too much time correcting the entry, the application will not be used regularly. The second was food-name specificity. A generic model may recognize a plate, bowl, or snack-like region, but it cannot reliably separate \code{rajma_curry}, \code{sambar}, \code{poha}, \code{thepla}, and \code{bhakarwadi}. For nutrition advice, those labels matter. A carbohydrate-heavy meal, for example, should not be explained in the same way for weight loss, muscle gain, and diabetes-related caution.

The connection between detection and nutrition lookup became clearer after the first working version was tested. A detector can correctly return \code{samosa} or \code{palak_paneer}, but the application still has to choose the right database row. If that mapping is wrong, the frontend can display neat-looking calorie and macro values that are actually misleading. One early fuzzy-matching result linked \code{palak_paneer} to a generic cottage-cheese entry instead of spinach paneer. That single case changed the design: nutrition mapping was moved into the main project workflow and was tested separately instead of being treated as a small helper function.

\section{Problem Statement}

The problem addressed in this project is to design and implement a full-stack system that can:

\begin{enumerate}[label=\arabic*.]
    \item Detect one or more Indian food items from an uploaded image.
    \item Normalize detected class labels into stable canonical food names.
    \item Map detected foods to nutrition database records with maximum safe coverage.
    \item Estimate calories, protein, carbohydrates, and fat per detected item.
    \item Provide personalized dietary recommendations using user goals and health conditions.
    \item Present the results through an accessible web interface with bounding boxes, serving adjustment, and macro tracking.
\end{enumerate}

The problem is therefore wider than training a detector. It includes dataset engineering, data validation, backend API design, frontend interaction design, and nutrition-aware application logic. Each part can affect the final result. A clean frontend cannot fix a wrong food-to-nutrition mapping, and a trained detector becomes less useful if the class labels are inconsistent.

\section{Objectives}

The main objectives of the project are:

\begin{itemize}
    \item To build a canonical Indian food dataset by merging multiple YOLO-format source datasets.
    \item To resolve duplicate and inconsistent food labels into a stable set of classes.
    \item To integrate a YOLO-based detection service into a FastAPI backend.
    \item To create a nutrition lookup service backed by a local SQLite database.
    \item To improve nutrition mapping quality through explicit class-to-database mappings.
    \item To implement macro calculation for different dietary goals and health profiles.
    \item To develop a Next.js frontend that supports image upload, detection visualization, nutrition cards, and recommendation display.
    \item To validate the dataset, class names, nutrition mappings, and service behavior.
\end{itemize}

\section{Scope of the Project}

The implemented prototype is limited to image-based recognition of Indian food classes present in the curated dataset. The model output is treated as a dish label, and nutrition values are reported per 100 grams unless the user changes the serving assumption in the frontend. The application does not estimate physical food volume or exact portion size from a single photograph. This was an intentional boundary. Reliable portion estimation would need depth information, a known reference scale, segmentation, or controlled image capture, none of which were available in the current prototype.

The backend and frontend are designed for local full-stack deployment. The backend exposes REST-style endpoints for image analysis, macro calculation, validation, and recommendation generation. The frontend is built as a single-page workflow for uploading images and reviewing results. When configured, the recommendation module uses Groq-backed language model generation, but the numerical nutrition calculations remain inside the backend. This separation keeps recommendations from becoming the source of calorie or macro values.

\section{Major Contributions}

The major contributions are listed below. They are written from the completed implementation rather than only from the model-training stage, because the final output depends on the full pipeline.

\begin{enumerate}[label=\arabic*.]
    \item A merged Indian food dataset containing 72 canonical classes and 48,693 exported image-label pairs.
    \item A normalization pipeline that maps raw labels such as \code{AlooGobi}, \code{aloo gobhi}, \code{thengai chutney}, \code{satham}, and \code{medu vadai} into stable canonical labels.
    \item Visual class verification artifacts showing real image samples for all 72 classes.
    \item A nutrition database build process that imports INDB nutrition fields and creates a runtime SQLite database.
    \item A safer nutrition mapping strategy that covers all 72 dataset classes using explicit semantic mappings and verified supplemental nutrition rows before fuzzy matching.
    \item A FastAPI analysis endpoint that combines detection results, bounding boxes, nutrition data, and food-not-found handling.
    \item A frontend interface that combines upload, detection visualization, serving adjustment, macro progress, and recommendations.
\end{enumerate}

\section{Report Organization}

The report follows the same broad order as the project work. Chapter 2 reviews the supplied research papers and identifies the gap addressed by the implementation. Chapter 3 explains dataset construction, nutrition mapping, macro logic, and validation strategy. Chapter 4 presents the system design and implementation details in one combined chapter. Chapter 5 presents the experimental analysis, results, and discussion together so that the training evidence and interpretation remain in one place. Chapters 6 and 7 contain the conclusion and future work. The appendices include class mappings, API examples, and supporting verification material.
